\section{Discussion}
\label{sec:discussion}

The proposed evaluation is designed to test the core thesis that agentic software engineering is more than a language-modeling task; it is a distributed systems problem that requires formal control planes.

\subsection{Systems Framing vs. Prompt Engineering}
Historically, the primary method for improving coding agent performance has been prompt engineering—tuning system instructions, adding few-shot examples, or structuring chat outputs. While these methods improve the probability of a model generating correct syntax, they do not change the fundamental execution model. A model executing directly on a production workspace remains unbounded and unverified.

Enforcing explicit custody and proof shifts the focus from inputs (prompts) to execution boundaries and evidence. Decapod can require a claimed task to enter an isolated workspace and can block promotion paths when validation or proof requirements are missing. This is narrower than complete sandboxing: the current kernel does not guarantee that every incorrect change is detected, that every credential is inaccessible, or that every agent action is intercepted. The evaluation must therefore measure both safety benefits and governance overhead.

\subsection{The Control Plane for Machine Work}
In traditional software engineering, human developers follow processes: they checkout branches, write tests, request pull request reviews, and run CI/CD pipelines. However, these processes were designed for human timelines. An autonomous agent can run dozens of commands per minute, coordinate across multiple repositories, and execute complex workflows concurrently.

Decapod acts as a local control plane for machine work. It brings intent resolution, task ownership, workspace isolation, validation, proof events, and publication gates into a repository-native workflow. It does not itself coordinate every command issued by an arbitrary model runtime, so the boundary between Decapod evidence and the agent host remains part of the threat model.

\subsection{Trust, Auditing, and the Limitations of Proof}
A central hypothesis of this study is that structured proof files and trajectories will reduce human review time. In a transcript-only workflow, human reviews are cognitively demanding: the developer must read long chat threads, try to understand why the agent chose a specific path, and then manually check if the code works. A Decapod run, however, separates the conversational noise from the systems evidence.

Programmatic proof has strict limits. Decapod evidence can show that configured commands ran and what they returned, but it cannot verify soft requirements such as code elegance, architectural maintainability, or completeness of an unstated user goal. The goal is therefore to make objective checks and governance state easier to inspect, not to replace human review or convert a validation pass into a security guarantee.
